\documentclass[10pt,a4paper]{article}
\usepackage[T1]{fontenc}
\usepackage{lmodern,amsmath,microtype,booktabs}
\usepackage[margin=20mm]{geometry}
\usepackage[hidelinks]{hyperref}
\pdfinfoomitdate=1
\pdftrailerid{}
\pdfsuppressptexinfo=15
\setlength{\parindent}{0pt}
\setlength{\parskip}{5pt}
\begin{document}
\begin{center}
{\Large Maintaining an extended, dim H$\alpha$ absorber}\par\vspace{4pt}
\texttt{exojax-escape} --- 27 November 2026
\end{center}

\textbf{Result.} The most concrete candidate is a denser, slowly moving outer
sector whose ionization/recombination balance continuously supplies excited
hydrogen, while non-LTE radiative rates keep its H$\alpha$ source low.
Magnetic or external-pressure confinement could supply the missing dynamics.
The local atomic experiment establishes conditional compatibility, not a
self-consistent atmosphere or an improved observational residual. Three
parallel investigations examine atomic budgets, irradiation and confinement;
production sources and the previous native-data comparisons are unchanged.

\textbf{What must be kept large, and what must stay small.}
For a resolved transition $l\to u$, under the common-profile convention,
\[
 \chi_\nu\propto n_l B_{lu}-n_u B_{ul},\qquad
 S_\nu=\frac{2h\nu^3}{c^2}
 \left(\frac{g_u n_l}{g_l n_u}-1\right)^{-1}.
\]
H$\alpha$ absorption mainly requires $n=2$ hydrogen. Its source depends on
the upper/lower population ratio, not on total density alone. Increasing both
populations can retain a small source while increasing opacity \emph{and}
absolute emissivity. The free-model clue is therefore not a non-emitting screen.
Above $1.725R_p$, original model 1 has $N_2=1.190\times10^{13}$ versus
$8.405\times10^{10}\,\mathrm{cm}^{-2}$. At $1.9375R_p$, it needs
$n_2=1401.5$ versus $13.14\,\mathrm{cm}^{-3}$, but its local nominal-center
source is only 1.323 times the physical source. Its inner emission and source
transition also differ; matching this outer clue does not reproduce model 1.

The saved outer gas is already far from LTE: at $7902$ K,
$n_3/n_2=0.00567$, versus the LTE value 0.14046. The three H$\alpha$ branch
excitation temperatures are $3614$--$3786$ K. Thus gas temperature is not the
temperature inferred from the line source. At the current irradiation,
$n=2$ photoionization alone acts in roughly $50$--$75\,\mu$s; individual
excited atoms cannot simply be carried outward and stored. This does not
measure the relaxation time of the gas's full ionization or thermal state.

\textbf{Recovered atomic balance and controlled response.}
The historical rate object omits later metal and H/O receiving corrections.
Restoring these for 2s, 2p, 3s, 3p and 3d reproduces their saved balance in
cells 4--6 to $5.2\times10^{-13}$ relative net/turnover. Ground-only chemical
terms and the below-threshold IR packets do not add direct terms to these
five equations. Large effective H$\alpha$ flows mostly recycle population:
after combining $n=2+3$, higher-level radiative exchange supplies a net
$1.113\times10^5\,\mathrm{cm}^{-3}\,\mathrm{s}^{-1}$ at cell 5; base
bound-free and ground radiative exchange remove $8.608\times10^4$ and
$1.821\times10^4$, respectively. The origin of the imposed higher-level
population is not identified by this decomposition.

A five-level solve scales electrons, protons and the external ground/high-level
baths by a density factor $D$, recomputes collisions and recombination, and
holds temperature, pumping and escape fixed:
\begin{center}
\begin{tabular}{rrrr}
\toprule
$r/R_p$ & Required $n_2$ factor & $D$ matching that $n_2$ & $S/S_0$ \\
\midrule
1.5125 & 4.72 & 3.15 & 0.976 \\
1.9375 & 106.63 & 20.19 & 0.978 \\
2.3625 & 323.15 & 36.56 & 0.981 \\
\bottomrule
\end{tabular}
\end{center}
These are independent local probes, not one solved radial profile. At cell 5,
$D=20.19$ yields $n_3=7.47\,\mathrm{cm}^{-3}$ and $n_3/n_2=0.00533$.
The fixed-bath forcing decomposition attributes 84.1\% of maintained $n_2$
to direct base recombination, versus 22.1\% at baseline; these are linear
response fractions, not independent gross/net source fractions.
Collisions do not immediately thermalize H$\alpha$ in this experiment.
The density factors are \emph{not inferred densities}: higher levels,
ionization, heating/cooling and the radiation field must all readjust.

\clearpage
\textbf{An immediate limiting feedback is radiative.}
The density probe raises local center opacity by 106.66, taking the cell's
static optical depth from about 0.041 to 4.35. Keeping its original escape
and pumping is then only a diagnostic. Multiplying both upward and downward
effective H$\alpha$ radiative rates by 0.25 leaves $n_2\simeq1376$ but
raises $S/S_0$ to 1.336; a factor 0.1 gives $n_2\simeq1332$ and $S/S_0=1.989$.
These are sensitivity probes, not recomputed escape probabilities. They show
that the excited reservoir can persist while its source becomes brighter.

\textbf{Why a slow, confined region merits priority.}
The saved flow has $\dot M=1.0513\times10^{13}\,\mathrm{g\,s^{-1}}$.
At cell 5, $v_r=4.879\,\mathrm{km\,s^{-1}}$ and
$P=2.228\times10^{-3}\,\mathrm{dyn\,cm^{-2}}$. At unchanged temperature
and composition, $D=20.19$ would require $v_r=0.242\,\mathrm{km\,s^{-1}}$
if the same spherical mass flux is retained, or 20.19 times more mass flux
at unchanged speed. A closed sector need not carry that global escaping flux.
A new force balance or boundary is necessary; slowing the prescribed
transonic solution by hand is not a mechanism.

For this illustrative density, $\sqrt{8\pi P}=1.06$ G locally. An ideal
equatorial dipole gives a surface field of 7.73 G (polar value twice as
large). This is a conditional gas-pressure scale, \emph{not} a field
measurement or a sufficient confinement criterion. Magnetic tension,
along-field forces, tides, open-field loading and the stellar environment
remain unsolved. Closed-field dense regions have precedents in hot-Jupiter
MHD calculations, but those Ly$\alpha$ transit studies do not establish this
KELT-9b H$\alpha$ emission mechanism
[\href{https://arxiv.org/abs/1404.5817}{Trammell et al. 2014}].

The present wind already includes $3\Omega_{\rm orb}^2r$ tides. With its
actual $2.17M_J$ model mass, the Hill/L1 distance is $2.474R_p$, but the
same critical equipotential intersects the current observer direction at
only $1.814$--$1.867R_p$. A spherical static shell at the free model's
$2.246R_p$ radial optical-depth radius cannot be justified merely because
it is inside L1. Three-dimensional support and coverage matter.

\textbf{Radiation mechanisms: useful coupling, weak stand-alone shortcuts.}
Ly$\alpha$ trapping, mixing, two-photon loss, Balmer photoionization and
recombination already exist. The Ly$\alpha$ escape probability at cell 5 is
$5.82\times10^{-6}$. Removing only its spontaneous escape term increases
$n_2$ by 7.65\% in the fixed-bath solve. Removing all $n=2,3$ photoionization
sinks increases it by 69.2\%, far short of the factor 107. These restricted
probes do not bound a new atmosphere whose other populations also respond.

Increasing only the effective Ly$\alpha$ pump needs a factor 3424 to reach
the target. It then produces $2.139\times10^7$ $n=2$ photoionizations per
$\mathrm{cm}^3$ per second, versus all 18 levels' base radiative recombination
of $3.726\times10^5$. The selected-level change in net proton production is
$+2.171\times10^7$ with other baths unchanged, explicitly breaking the saved
ionization balance. Absolute full-atom receiving rates were not recovered,
but their unchanged contributions cancel in this fixed-bath difference.
More pump alone is therefore not a closed solution.

Nonlocal Ly$\alpha$ transport from inner shells, angle dependence and partial
redistribution could change where excitation is maintained. The current
escape closure does not solve that spatial/frequency problem. This is a
closure test, not a missing Ly$\alpha$ reaction; redistribution need not
increase trapping [\href{https://arxiv.org/html/2304.07352v1#S2.SS7}{Huang et al. 2023}].

\clearpage
\textbf{A UV shield must absorb the right photons.}
Ground-state H absorbs above 13.6 eV, whereas $n=2$ is photoionized above
3.4 eV. In the adopted unattenuated stellar spectrum, 88--93\% of its
photoionization is supplied by 3.4--6 eV photons. The outer ground-H column
has optical depth 7.16 at 13.7 eV, yet no ground-H bound-free opacity below
threshold. An EUV shield is not a Balmer-continuum shield.
At 3.41 eV, the current total outer continuum depth is only
$1.156\times10^{-6}$; the free target $N_2$ gives an H($n=2$) depth of
$1.635\times10^{-4}$. Even at this favorable energy, a pure $n=2$ screen
attenuating by 142 requires $N_2\simeq3.60\times10^{17}\,\mathrm{cm}^{-2}$,
about 30,000 times the target. This is a monochromatic attenuation condition,
not a prediction of a 142-fold population increase.

Additional excited-metal blanketing or an external dense screen is not
excluded, but needs broad-band opacity, surviving species and consistent
re-emission/heating. Balmer absorption already heats the atmosphere;
suppressing it changes recombination and flow as well as the $n=2$ sink
[\href{https://arxiv.org/html/1910.00267v1}{Garc\'ia Mu\~noz \& Schneider 2019}].
At cell 5, saved base-field $n=3$ photoionization is only 0.2--4\% of the bound-bound
radiative outflow rates, so modest optical/Paschen changes are not an obvious
way to strongly lower $n_3$. A scalar finite-stellar-disk correction changes
the mean intensity by only 2.63\%. Illustrative stellar velocity averaging
actually raises H$\alpha$ pumping by factors 1.39--1.60 for the supplied deep
stellar line. These kernels do not reconstruct the planet-view stellar spectrum.

\textbf{Dim local gas need not make a dim unresolved spectrum.}
Along a covered ray,
\[
 I_\nu=I_{\nu,\rm inner}e^{-\tau_\nu}
        +S_{\nu,\rm outer}(1-e^{-\tau_\nu}).
\]
Off the inner disk, the same gas emits. In a uniform opaque-disk proxy,
$F_\nu\propto R_{\rm eff}^2S_\nu$. Substituting the two profiles' static
radial $\tau=1$ radii gives an area factor 1.991 and a source factor 0.605:
their product is 1.204, \emph{larger} despite the lower source. These radii
are not observer formation radii and this is not a rendered flux prediction;
it is a counterexample to an automatic dimming inference. A fully opaque
shell around an emitter of radius $1.3R_p$ and extending to $2.246R_p$
would need $S/I_{\rm inner}<0.335$ in this simplified limit to reduce the
total flux relative to the bare emitter.

A restricted foreground sector could improve this area tradeoff, but its
covering fraction, velocities and re-emission must be calculated. At the
selected phases the dayside is viewed at roughly $38$--$45$ degrees from
the star-facing axis; an anti-stellar planetary shadow is not automatically
in front of that emitter. Clumps likewise require thermal/pressure support
and projected covering, not just a recombination multiplier.

\textbf{Next discriminating calculation.}
Start with a specified slow/dense sector and the ordinary escaping region.
Re-solve the full populations, charge and energy balance with the existing
continuum channels; compare the current Ly$\alpha$ closure with spatial
transport, and solve H$\alpha$ pumping/escape at the new opacity. Retain only
states admitting the assumed force/boundary support. Render the disk and
halo together through the same native observation operator, including the
remaining phases and red wing. This can first test atomic/radiative viability
without committing to a complete MHD implementation. A dense sector is a
ranked hypothesis; no new fit, calibrated detection or production readiness
is claimed here.

\textbf{Evidence and verification.}
The campaign under
\texttt{mceg/evidence/halpha\_outer\_reservoir\_20261127/}
contains replay helpers, input/source pins, full local budgets, independent
reviews and primary-source scope notes. MCEG records finite calculation
support separately from the unknown self-consistent mechanism. The existing
phase-localization uncertainty and the preceding Decision history remain.
\end{document}
